\documentclass[../main.tex]{subfiles}
\begin{document}

Chương 6 tổng kết các kết quả chính của luận văn, thảo luận những hạn chế còn tồn đọng
và đề xuất hướng phát triển tiếp theo. Nội dung chương được tổ chức theo hai phần tương
ứng với hai câu hỏi trung tâm: đề xuất thiết kế Guepard Shield đã được kiểm chứng ở mức
độ nào, và những điều kiện nào cần được đáp ứng để hệ thống sẵn sàng cho triển khai thực
tế.

\section{Tổng kết}
\label{sec:conclusions_summary}

Luận văn đề xuất và xây dựng Guepard Shield, một pipeline đầu-cuối cho phát hiện bất
thường dựa trên syscall theo hướng lai giữa học sâu và thực thi nhẹ trong nhân hệ điều
hành. Xuất phát điểm của nghiên cứu là khoảng cách thực thi chính sách bảo mật (enforcement
gap): các mô hình Transformer có năng lực học mẫu tuần tự tốt nhưng không thể chạy trong
đường đi syscall của kernel, trong khi các hệ thống dựa trên luật tĩnh như Falco hay
Tetragon có chi phí thấp nhưng phụ thuộc vào tri thức chuyên gia thủ công. Guepard Shield
lấp khoảng trống đó bằng cách tách giai đoạn học ngoại tuyến ra khỏi giai đoạn thực thi
thời gian thực, sử dụng mô hình Transformer như Teacher để học phân phối hành vi bình
thường, sau đó chưng cất kết quả học được thành một DFA Student đủ nhỏ để nạp vào BPF map.

\subsection{Các kết quả đã đạt được}

Luận văn xác nhận tính khả thi của pipeline qua bốn nhóm kết quả chính.

\textbf{Thiết kế pipeline đầu-cuối.} Ba giai đoạn từ tiền xử lý dữ liệu (Giai đoạn 1),
huấn luyện Transformer Teacher (Giai đoạn 2), đến chưng cất DFA Student qua K-Means (Giai
đoạn 3) đã được xây dựng và kết nối thành một quy trình liên tục. Ranh giới rõ ràng giữa
offline learning và online enforcement là nguyên tắc thiết kế quan trọng giúp mỗi giai đoạn
có thể được tối ưu độc lập mà không làm ảnh hưởng đến chi phí runtime.

\textbf{Transformer Teacher trên LID-DS-2021.} Mô hình Decoder-only Transformer với
$d_{\text{model}} = 128$, 4 lớp và 4 attention heads, huấn luyện theo mục tiêu dự đoán
token kế tiếp trên dữ liệu bình thường, đạt validation loss 0.3686 tại epoch 1 và hội tụ
ổn định qua 7 epoch. Điểm last-token NLL đạt AUROC = 0.8503 và PR-AUC = 0.7093 trên gần
100 triệu window kiểm thử, cho thấy mô hình có khả năng phân tách window bất thường khỏi
hành vi bình thường ngay cả khi chỉ học từ dữ liệu không giám sát. Ở cấp độ bản ghi, 1,715
trong 1,815 attack recordings (94.5\%) có ít nhất một window vượt ngưỡng oracle, với
median detection lag là 14 window sau mốc khai thác.

\textbf{Chưng cất DFA Student.} Cấu hình S3 Majority Voting tại $K = 64$ tạo ra bảng
chuyển trạng thái với 64 states và 1,973 transitions, đạt tỷ lệ phát hiện cấp độ bản ghi
$DR_{\text{rec}} = 90.85\%$, tỷ lệ báo động giả FPR = 1.41\% và fidelity 98.12\% so với
Teacher. Kích thước bảng chuyển dưới dạng mảng trực tiếp $64 \times 816 \times 4$ bytes
xấp xỉ 204 KB, hoàn toàn nằm trong giới hạn BPF map của các kernel Linux hiện đại.
Phân tích đánh đổi giữa các chiến lược giải quyết tính không xác định cho thấy majority
voting là lựa chọn khả dụng duy nhất; statistical pruning S4 thất bại ở mọi giá trị $K$
được thử nghiệm do tỷ lệ non-determinism không đủ thấp để pruning giữ lại bảng chuyển có
ý nghĩa.

\textbf{Phân tích lý thuyết và khung đánh giá.} Chương 4 xác lập độ phức tạp $O(1)$ cho
mỗi sự kiện syscall tại runtime, phân tích giới hạn bộ nhớ BPF map, khung đánh giá fidelity
và đặc tính kháng mimicry attacks của DFA whitelist. Khung đánh giá này cũng làm rõ rằng
nhiễu nhãn cấu trúc trong LID-DS-2021, nơi phần lớn tail sau khai thác vẫn mang hành vi
bình thường, là lý do tại sao TPR cấp độ window thấp không mâu thuẫn với $DR_{\text{rec}}$
cao: một cơ chế giám sát đúng nên chấp nhận post-exploit normal behavior thay vì cố gắng
gắn cảnh báo cho toàn bộ phần tail được gán nhãn attack.

\subsection{Các vấn đề còn tồn đọng}

Bên cạnh các kết quả trên, luận văn thừa nhận một số vấn đề chưa được giải quyết đầy đủ.

Thứ nhất, thực nghiệm trên miền $K \in \{128, 256, 512\}$ xác nhận
rằng các cấu hình này không khả dụng và không phù hợp để phát triển
tiếp. Tỷ lệ empty clusters cao tại các giá trị $K$ lớn hơn cho thấy
không gian hidden state của Teacher có số chiều hiệu dụng thấp hơn
đáng kể so với những giá trị $K$ đó; việc ép phân cụm vào nhiều
cluster hơn chỉ làm tăng non-determinism và FPR mà không cải thiện
fidelity hay $DR_{\text{rec}}$. Chiến lược S4 statistical pruning vì
vậy cũng bị loại khỏi xem xét, vì điều kiện tiên quyết của nó là tỷ
lệ non-determinism thấp không được đáp ứng ở bất kỳ $K$ nào được
thử nghiệm.

Thứ hai, phân tích MITRE ATT\&CK mới dừng lại ở mức khảo sát khả năng ánh xạ. Việc gắn
nhãn các trạng thái attack trong DFA với các kỹ thuật tấn công cụ thể như leo thang đặc
quyền hay chiếm quyền điều khiển tiến trình đòi hỏi phân tích sâu hơn về các recording
tấn công trong LID-DS và bộ dữ liệu DongTing.

Thứ ba, mô hình Teacher trong luận văn chỉ sử dụng tên syscall như đặc trưng đầu vào,
chưa khai thác tham số syscall (như đường dẫn tệp trong \texttt{open}, địa chỉ trong
\texttt{connect}) hay siêu dữ liệu tiến trình. Thiếu thông tin này làm giảm khả năng phân
biệt các hành vi có cùng chuỗi syscall nhưng khác nhau về tham số.

\section{Hướng phát triển trong tương lai}
\label{sec:future_works}

Dựa trên các kết quả và hạn chế đã trình bày, luận văn đề xuất hai hướng phát triển tiếp
theo.

\textbf{Mở rộng biểu diễn đầu vào bằng tham số syscall và siêu dữ liệu tiến trình.} Giới
hạn lớn nhất của cách token hóa hiện tại là chỉ dùng tên syscall. Một hướng cải thiện
khả thi là xây dựng từ vựng tổng hợp từ cặp (syscall, nhóm tham số), ví dụ phân biệt
\texttt{open}/đường dẫn nội bộ với \texttt{open}/\texttt{/proc/}, hoặc bổ sung các trường
siêu dữ liệu như UID, GID và tên tiến trình cha vào ngữ cảnh đầu vào. Cách tiếp cận này
làm tăng kích thước từ vựng $|\Sigma|$ và có thể làm tăng không gian DFA, nhưng có thể
cải thiện đáng kể khả năng phân biệt hành vi khai thác với hành vi bình thường có cùng
chuỗi syscall.

\textbf{Đánh giá trên bộ dữ liệu đa dạng và mở rộng ánh xạ MITRE ATT\&CK.} LID-DS-2021
là bộ dữ liệu đánh giá chính nhưng được thu thập trong môi trường container hóa tương đối
đồng nhất. Để đánh giá khả năng tổng quát hóa của pipeline, các thực nghiệm tiếp theo
nên bao gồm bộ dữ liệu DongTing thu thập từ hệ thống thực tế, cũng như các kịch bản
fileless attacks và memory exploitation không để lại dấu vết trên đĩa. Song song với đó,
việc xây dựng ánh xạ có hệ thống giữa các trạng thái DFA bị gắn attack với các kỹ thuật
tấn công trong ma trận MITRE ATT\&CK sẽ giúp các analyst có thêm ngữ cảnh diễn giải khi
nhận cảnh báo từ hệ thống, thu hẹp khoảng cách giữa phát hiện bất thường tự động và hiểu
biết về ý đồ của kẻ tấn công.

\end{document}
